% Generated draft: requires mathematical and source review.
\documentclass[11pt]{article}
\usepackage[margin=25mm]{geometry}
\usepackage{fontspec}
\setmainfont{Latin Modern Roman}
\usepackage{amsmath,amssymb,mathtools,unicode-math}
\setmathfont{Latin Modern Math}
\usepackage{enumitem,needspace,xurl,hyperref}
\setlength{\parindent}{0pt}
\setlength{\parskip}{5pt}
\setlength{\emergencystretch}{3em}
\newcommand{\N}{\mathbb{N}}
\newcommand{\Z}{\mathbb{Z}}
\newcommand{\Q}{\mathbb{Q}}
\newcommand{\R}{\mathbb{R}}
\newcommand{\C}{\mathbb{C}}
\newcommand{\F}{\mathbb{F}}
\newcommand{\A}{\mathbb{A}}
\newcommand{\PP}{\mathbb{P}}
\newcommand{\E}{\mathbb{E}}
\newcommand{\eps}{\varepsilon}
\newcommand{\dd}{\,\mathrm{d}}
\begin{document}
\section*{UnsolvedMath ID: 30002760}
\textbf{Source problem code:} catalogue entry
\subsection{Definitions and mathematical statement}
The supplied problem concerns adaptive finite-element approximation of elliptic or parabolic transport-dominated diffusion equations. Its background specifically mentions isotropic refinement classes generated by newest-vertex bisection; the statement also allows a refinement method ``such as'' newest-vertex bisection. The equation, domain, coefficients, data, boundary conditions, and, in the parabolic case, initial conditions and time interval are not specified.

A conditional formulation of the requested rate comparison is as follows. Once a particular equation with solution \(u\), a finite-element method, an error measure, a refinement class \(\mathbb T\), and a positive integer complexity measure \(n(\mathcal T)\) have been specified, let \(e(u,\mathcal T)\geq 0\) denote the error of that method's approximation on \(\mathcal T\). Define the best partition error at complexity budget \(N\) by
\[
\sigma_N(u):=\inf\bigl\{e(u,\mathcal T):\mathcal T\in\mathbb T,\ n(\mathcal T)\leq N\bigr\},\qquad N\in\N,\ N\geq1,
\]
with \(\inf\varnothing=+\infty\). If an adaptive algorithm produces meshes \((\mathcal T_k)_{k\geq0}\subseteq\mathbb T\), write \(n_k=n(\mathcal T_k)\).

\emph{Rate-optimality question, conditional on those choices.} Does the adaptive approximation attain every algebraic convergence rate attainable by the best admissible partitions? A schematic expression of this question is whether, for every \(s>0\),
\[
\begin{aligned}
&\bigl[\exists A>0\ \exists N_0\in\N\ \forall N\geq\max\{1,N_0\}:\ \sigma_N(u)\leq A N^{-s}\bigr]\\
&\qquad\Longrightarrow
\bigl[\exists B>0\ \exists k_0\in\N\ \forall k\geq k_0:\ e(u,\mathcal T_k)\leq B n_k^{-s}\bigr],
\end{aligned}
\]
for adaptive sequences whose complexities tend to infinity.

This display is a conditional interpretation of convergence-rate optimality, not a fully specified conjecture supplied by the record. The record does not determine the error, complexity measure, admissible equation class, adaptive algorithm, dependence of constants, or whether only algebraic rates are intended. It also does not determine whether the best partition error uses computed finite-element solutions or best approximation within the finite-element spaces. In the parabolic case, the treatment of time discretization and its cost is unspecified. These choices must be resolved before a precise theorem target can be fixed.
\subsection{Short English statement}
For elliptic or parabolic transport-dominated diffusion equations, does adaptive finite-element approximation converge as quickly, relative to partition complexity, as the best partitions obtainable by the allowed refinement procedure? The supplied background emphasizes isotropic newest-vertex bisection. Essential choices of equation, error measure, algorithm, and complexity are missing, so the mathematical display is a conditional rate-comparison template requiring review.
\subsection{Sources}
\begin{itemize}
\item Local UnsolvedMath registry, supplied source record: permanent ID \(30002760\), title ``Optimal Adaptive Approximation of Transport-Dominated Equations,'' source problem code \texttt{catalogue entry}, in the Oberwolfach Reports Open Problems collection. The recorded proposed year is 2015; no individual proposer is supplied.
\item Supplied report citation: \emph{New Discretization Methods for the Numerical Approximation of PDEs} (2015), Oberwolfach Reports, pp.~87--185, DOI \texttt{10.4171/owr/2015/2}. \url{https://doi.org/10.4171/owr/2015/2}. The record describes Oberwolfach Reports as documenting mathematical workshops and their open-problem discussions; it does not supply a problem-specific page or passage.
\item Supplied literature source, labeled ``original OWR report'': \url{https://doi.org/10.4171/owr/2015/2}.
\item Supplied literature source, labeled ``current report summary of remaining challenges'': \url{https://oa.tib.eu/renate/items/3429c126-b128-463f-9075-31c40105b9ef/full}.
\item Literature assessment dated exactly \texttt{2026-08-21}: the supplied curation labels the problem open and states that, despite advances in adaptive finite-element theory, the report's exact optimality question for transport-dominated diffusion and isotropic newest-vertex-bisection classes remains an open challenge. This is supplied dated curation, not an independently verified assessment; its time-sensitive status and cited claims require independent verification.
\item The registry labels the statement \texttt{corrected\_verified} and supplies a claim of manual semantic, typographic, and LaTeX review confirming a self-contained, meaning-preserving cleanup. These are supplied review metadata, not independent verification of the report attribution, mathematical formulation, or present open status.
\end{itemize}
\end{document}
